\begin{table}[htbp]\centering\small\singlespacing
\caption{Country-level AERI: descriptive comparisons, 2015--2025}\label{tab:country_aeri}
\begin{tabular}{lrrrrl}\toprule
Country & Months & Mean & SD & Peak & Peak month\\\midrule
Nigeria & 132 & 73.04 & 8.75 & 98.43 & 2024-02\\
Ghana & 132 & 57.11 & 10.47 & 85.70 & 2025-03\\
Algeria & 132 & 65.02 & 6.34 & 83.44 & 2020-04\\
Morocco & 132 & 22.67 & 7.54 & 37.66 & 2025-10\\
Egypt & 132 & 40.55 & 11.29 & 97.11 & 2025-12\\
Ethiopia & 132 & 59.96 & 17.50 & 127.83 & 2019-02\\
South Africa & 132 & 42.58 & 7.61 & 79.22 & 2020-04\\
Kenya & 132 & 38.85 & 7.02 & 71.03 & 2025-12\\
Angola & 132 & 53.45 & 8.60 & 85.67 & 2019-10\\
Ivory Coast & 132 & 30.22 & 5.28 & 46.87 & 2025-04\\
\bottomrule\end{tabular}
\par\medskip\raggedright\footnotesize\textit{Notes:} Monthly AERI uses the article-score definition reported in the paper. Only months with at least 100 retained articles are included. Means and sample standard deviations give equal weight to eligible months; they differ from article-weighted annual AERI. Peak denotes the maximum eligible monthly observation, with the earliest month reported in a tie. These ten economies form the illustrative group used in the paper; PVAR sample selection is conducted separately.
\end{table}